\documentclass[11pt]{article}
\usepackage[margin=1in]{geometry}
\usepackage{amsmath,amssymb,amsthm,enumitem,hyperref}
\hypersetup{colorlinks=true,linkcolor=blue,urlcolor=blue,citecolor=blue}
\newtheorem{theorem}{Theorem}
\newtheorem{lemma}{Lemma}
\newcommand{\Q}{\mathcal Q}
\newcommand{\conv}{\operatorname{conv}}
\newcommand{\supp}{\operatorname{supp}}
\title{Simplicial polytopes realizing arbor support vectors\\\large MathDB 373576}
\author{Principia Math In Progress}
\date{September 18, 2026}
\begin{document}
\maketitle
\begin{abstract}
We prove Chapoton's conjecture that the support-count vector of every
arbor polytope is the $h$-vector of a simplicial polytope. In fact the
result holds for every finite preorder. We identify a special simplex
in an augmented bipartite root polytope and verify unimodularity.
The proof uses, with explicit attribution, the support-enumerator
theorem of Dai--Hou--Liu--Thawinrak--Wang and Athanasiadis's
special-simplex theorem.
\end{abstract}

\section*{Task 1. Full solution}
Let $\rho$ be a preorder on $E=[n]$, where $n\ge1$, and put
\[
 \Q_\rho=\{x\in\mathbb R_{\ge0}^{E}:x(I)\le |I|
                   \text{ for every order ideal $I$ of $\rho$}\},
 \qquad x(A)=\sum_{i\in A}x_i.
\]
An order ideal contains every predecessor of each of its elements;
equivalent elements of a preorder are therefore included together.
Write $P_\rho=\Q_\rho\cap\mathbb Z^E$ and
\[
 H_\rho(t)=\sum_{a\in P_\rho}t^{|\supp(a)|},
 \qquad \supp(a)=\{i:a_i>0\}.
\]
For a simplicial $n$-polytope $S$, the convention is
$h(\partial S,t)=\sum_{i=0}^n f_{i-1}(\partial S)t^i(1-t)^{n-i}$,
with $f_{-1}=1$.
\begin{theorem}\label{thm:main}
For every preorder $\rho$ on $[n]$ there is an $n$-dimensional
simplicial convex polytope $S_\rho$ such that
\[
 \boxed{H_\rho(t)=h(\partial S_\rho,t).}
\]
In particular this holds for the support-count polynomial of every
arbor, proving Conjecture 1.1 of \cite{AXY} in full. Its coefficients
are palindromic and unimodal, and satisfy the $g$-theorem.
\end{theorem}

\subsection*{A root polytope with the required Ehrhart numerator}
Form a bipartite graph $G$ with both vertex sets indexed by $E$,
placing an edge from left $i$ to right $j$ precisely when
$j\le_\rho i$. If $A$ is a set of left vertices, its neighborhood
is the ideal $\downarrow_\rho A$. Thus the set of draconian vectors
in the convention of \cite{D} is
\[
 D(G)=\{a\in\mathbb Z_{\ge0}^{E}:
                   a(A)\le |\downarrow_\rho A|\text{ for all }A\subseteq E\}
      =P_\rho.                                                    \tag{1}
\]
Indeed the inequalities for ideals give one inclusion; for the other,
$a(A)\le a(\downarrow A)\le|\downarrow A|$.

Adjoin left and right vertices indexed by $0$, each adjacent to every
vertex on the other side. In the augmented root polytope, suppressing
the two coordinates indexed by $0$ gives the integral affine model
\[
 R_\rho=\conv\bigl(0,(e_i,0),(0,e_j),(e_i,e_j):j\le_\rho i\bigr)
                  \ \subset\mathbb R^{E}\oplus\mathbb R^{E}.       \tag{2}
\]
The original coordinates have left sum and right sum both equal to
$1$, so the suppressed coordinates are $1-\sum\alpha_i$ and
$1-\sum\beta_j$. Hence this model preserves the affine lattice.
It has dimension $2n$, since $0,(e_i,0),(0,e_i)$ are affinely independent.
The established support-enumerator theorem \cite[Theorem 1.2]{D}
and (1) give
\[
 \sum_{k\ge0}|kR_\rho\cap\mathbb Z^{2n}|t^k
       =\frac{H_\rho(t)}{(1-t)^{2n+1}}.                             \tag{3}
\]
This previously proved identity is an input, not a new claim here.

\subsection*{Facet description and the special simplex}
We claim that (2) has the following inequality description:
\[
 \alpha_i,\beta_i\ge0,\qquad
 \beta(F)+\alpha(E\setminus F)\le1
                  \quad\text{for every filter $F$ of $\rho$}.       \tag{4}
\]
The empty and full filters are included. A filter contains all
successors of each of its elements.

Here are both directions, including sufficiency. Every vertex in (2)
satisfies (4): an edge with $j\in F$ has $i\in F$, since $j\le_\rho i$.
Conversely suppose (4) holds. Set
$\alpha_0=1-\alpha(E)$ and $\beta_0=1-\beta(E)$; these are nonnegative.
We seek a nonnegative edge weighting of the augmented graph with these
left and right marginals. The fractional Hall criterion (equivalently,
the max-flow/min-cut theorem) says that it exists exactly when, for
every set $B$ of right vertices, its demanded mass does not exceed the
mass of its left neighborhood. Sets containing $0$ have all left
vertices as neighbors, so their inequalities are automatic.
For $B\subseteq E$, the left neighborhood is
$\{0\}\cup\uparrow_\rho B$. Write $F=\uparrow_\rho B$. Then (4) gives
\[
 \beta(B)\le\beta(F)\le1-\alpha(E\setminus F)
                         =\alpha_0+\alpha(F).
\]
These are all the required cut inequalities. The resulting edge
weights sum to $1$ and express $(\alpha,\beta)$ as a convex combination
of (2), proving the claim.

Consider the $n$-simplex
\[
 \Sigma=\conv\{0,(e_1,e_1),\ldots,(e_n,e_n)\}\subset R_\rho.        \tag{5}
\]
All of these points are vertices: they are distinct $0$--$1$ vectors
in the vertex set of (2), and a $0$--$1$ vector cannot be a convex
combination of different $0$--$1$ vectors. They are affinely independent.
Every facet of $R_\rho$ is supported by one of the inequalities in (4),
because (4) is a finite full-dimensional description.
A coordinate facet $\alpha_i=0$ or $\beta_i=0$ contains exactly the
$n$ vertices of $\Sigma$ other than $(e_i,e_i)$.
A filter facet contains every $(e_i,e_i)$ and omits $0$.
Consequently every facet contains exactly $n$ of the $n+1$ vertices
in (5). This is precisely the special-simplex condition.

\subsection*{Unimodularity and polytopality}
Every full-dimensional simplex with vertices from (2) is unimodular.
To see this, return to the augmented bipartite graph. A set of its
edge vectors is affinely independent exactly when the edges form a
forest. A cycle gives the alternating affine dependence; conversely,
successively removing leaves proves independence for a forest.
A full simplex has $2n+1$ edges on $2n+2$ graph vertices, so these
edges form a spanning tree.

For any integral left and right marginals with the same total, the
unique real edge weights on such a tree are integral: the weight on
an edge incident to a leaf is the marginal at that leaf; subtract it
from the adjacent marginal and continue. Nonnegativity is irrelevant
to this lattice assertion. Thus the tree edge vectors form an affine
lattice basis in total degree $1$, proving normalized volume one.
It follows that every triangulation using the vertices of $R_\rho$
is unimodular. In particular $R_\rho$ is compressed: every pulling
triangulation is unimodular.

We now apply the established special-simplex theorem
\cite[Theorem 3.5]{A}: if a $d$-dimensional integral polytope has a
unimodular pulling triangulation in which the $q$ vertices of a
special simplex are pulled first, then its Ehrhart numerator is the
$h$-polynomial of a simplicial polytope of dimension $d-q+1$.
All its hypotheses have just been verified, with $d=2n$ and $q=n+1$.
In combination with (3), its conclusion is exactly Theorem~\ref{thm:main}.

More concretely, the pulling triangulation is the join of $\Sigma$
with a complex $\Delta$ that triangulates the boundary of the quotient
by the diagonal subspace. This quotient is
\[
 A_\rho=\conv\{\pm e_i,\ e_j-e_i:j\le_\rho i\}\subset\mathbb R^E.
                                                                    \tag{6}
\]
Theorem 3.5 realizes $\Delta$ as the boundary of a simplicial polytope
by successive vertex pulling. Since the simplex factor has
$h$-polynomial $1$, $h(\Delta,t)=h^*(R_\rho,t)=H_\rho(t)$.
The quotient (6) itself need not already be simplicial.

\subsection*{The arbor target and the precise scope}
An arbor is a rooted tree of nonempty blocks partitioning $[n]$.
Give its labels the preorder $i\le_\rho j$ when the block of $i$
is a descendant of the block of $j$. For each tree vertex $v$, let
$D(v)$ be the union of the blocks of its descendants, including $v$.
Every ideal of $\rho$ is the disjoint union of the sets $D(v)$
at its maximal blocks. Therefore the arbor inequalities
$x(D(v))\le|D(v)|$ are equivalent to all the ideal inequalities
defining $\Q_\rho$. Its support-count vector is exactly the vector
in the question, with no restriction on branching or block sizes.

Palindromicity and unimodality now follow from the $g$-theorem for
simplicial polytopes. The stronger flag-realization and
gamma-positivity conjectures are not established by this argument.
The linear-arbor flag result of \cite{linear} is earlier work.
The support-root-polytope identity and the special-simplex theorem
are also earlier work; the contribution is their verified application
to the graph and special simplex above.

\section*{Tasks 2--5. Proposed follow-ups}
\begin{enumerate}[label=\textbf{Task \arabic*.},start=2,leftmargin=*]
\item Determine which arbors admit an ordering in the construction
above for which the resulting simplicial boundary is flag; in
particular, decide whether every arbor admits one.
\item Describe the facets of the quotient (6) directly from the
blocks and branches of an arbor, and give a combinatorial rule for
the maximal simplices of its boundary pulling triangulation.
\item Prove or disprove that the gamma-coefficients of $H_\rho(t)$
are nonnegative for every finite preorder $\rho$.
\item Construct a direct bijection from $P_\rho$ to the maximal
simplices of the special-simplex link, carrying support size to
restriction-face size in an explicit shelling.
\end{enumerate}

\paragraph{Verification.}
The exact verifier checks all 389 labeled preorders through size four,
including 12,708 unimodular maximal simplices and the complete
$h$-vectors of their links. Independent edge-sum and inequality counts
give 226 Ehrhart values for all 34 preorders through size three.
The arbor reduction is checked on all 1,674 labeled arbors through
size five, with 39 larger geometric examples. These finite checks
support the proof; they do not replace it. Review is by the authoring
agent only.

\begin{thebibliography}{9}
\bibitem{AXY} C. A. Athanasiadis, Q. Xiao, and X. Yan,
\emph{Lattice point enumeration of some arbor polytopes},
\href{https://arxiv.org/abs/2603.11654}{arXiv:2603.11654v1}, 2026,
Conjecture 1.1.
\bibitem{D} Z. Dai, Q. Hou, Z. Liu, W. Thawinrak, and H. Wang,
\emph{Counting Lattice Points in Minkowski Sums of Cross Polytopes},
\href{https://arxiv.org/abs/2608.16037}{arXiv:2608.16037v2}, 2026,
Theorem 1.2 and Section 4.2.
\bibitem{A} C. A. Athanasiadis,
\emph{Ehrhart polynomials, simplicial polytopes, magic squares and
a conjecture of Stanley}, J. Reine Angew. Math. \textbf{583} (2005),
163--174; \href{https://arxiv.org/abs/math/0312031}{arXiv:math/0312031},
Theorem 3.5.
\bibitem{linear} C. A. Athanasiadis,
\emph{Lattice Point Enumeration of Polytopes Associated to Integer
Compositions}, Ann. Comb. (2026),
\href{https://doi.org/10.1007/s00026-026-00812-2}{doi:10.1007/s00026-026-00812-2}.
\end{thebibliography}
\end{document}
